
\subsection{Identificación de recursos}
La identificación de recursos permite reconocer los medios humanos, metodológicos, tecnológicos e informacionales necesarios para abordar de manera adecuada la problemática del proyecto. Este apartado no solo cumple una función descriptiva, sino que también contribuye a evaluar la factibilidad del proceso de trabajo y a explicitar los apoyos requeridos para desarrollar la propuesta con coherencia.

\begin{itemize}
    \item \textbf{Recursos humanos:} equipo multidisciplinario del proyecto compuesto por estudiantes de Ingeniería Civil Telemática e Ingeniería en Diseño de Productos, docentes guía de ambas disciplinas, profesionales médicos del área asistencial consultados como informantes clave y futuros usuarios del sistema considerados para instancias de validación.

    \item \textbf{Recursos informacionales:} antecedentes sobre el traspaso clínico y cambio de turno en contextos hospitalarios, marcos normativos chilenos aplicables al registro clínico y la privacidad del dato, literatura sobre interoperabilidad de sistemas de salud y experiencia de usuario en entornos de alta carga cognitiva.

    \item \textbf{Recursos metodológicos:} herramientas de análisis sistémico, mapa de actores, buyer persona, mapa de empatía, modelación de la situación actual AS-IS y metodología Doble Diamante.

    \item \textbf{Recursos tecnológicos:} herramientas de diagramación y documentación, software de prototipado digital, tecnologías de captura de audio y procesamiento de lenguaje natural para la implementación del sistema y herramientas de análisis de flujos de interacción.
\end{itemize}
